 \documentclass[a4paper,11pt,oneside]{memoir}

% Castellano
\usepackage[spanish,es-tabla]{babel}
\selectlanguage{spanish}
\usepackage[utf8]{inputenc}
\usepackage{placeins}
\usepackage{eurosym}
\usepackage{float}

\RequirePackage{booktabs}
\RequirePackage[table]{xcolor}
\RequirePackage{xtab}
\RequirePackage{multirow}

% Links
\usepackage[colorlinks]{hyperref}
\hypersetup{
	allcolors = {blue}
}

% Ecuaciones
\usepackage{amsmath}

% Rutas de fichero / paquete
\newcommand{\ruta}[1]{{\sffamily #1}}

% Párrafos
\nonzeroparskip


% Imagenes
\usepackage{graphicx}
\newcommand{\imagen}[2]{
	\begin{figure}[!h]
		\centering
		\includegraphics[width=0.9\textwidth]{#1}
		\caption{#2}\label{fig:#1}
	\end{figure}
	\FloatBarrier
}

\newcommand{\imagenflotante}[2]{
	\begin{figure}%[!h]
		\centering
		\includegraphics[width=0.5\textwidth]{#1}
		\caption{#2}\label{fig:#1}
	\end{figure}
}



% El comando \figura nos permite insertar figuras comodamente, y utilizando
% siempre el mismo formato. Los parametros son:
% 1 -> Porcentaje del ancho de página que ocupará la figura (de 0 a 1)
% 2 --> Fichero de la imagen
% 3 --> Texto a pie de imagen
% 4 --> Etiqueta (label) para referencias
% 5 --> Opciones que queramos pasarle al \includegraphics
% 6 --> Opciones de posicionamiento a pasarle a \begin{figure}
\newcommand{\figuraConPosicion}[6]{%
  \setlength{\anchoFloat}{#1\textwidth}%
  \addtolength{\anchoFloat}{-4\fboxsep}%
  \setlength{\anchoFigura}{\anchoFloat}%
  \begin{figure}[#6]
    \begin{center}%
      \Ovalbox{%
        \begin{minipage}{\anchoFloat}%
          \begin{center}%
            \includegraphics[width=\anchoFigura,#5]{#2}%
            \caption{#3}%
            \label{#4}%
          \end{center}%
        \end{minipage}
      }%
    \end{center}%
  \end{figure}%
}

%
% Comando para incluir imágenes en formato apaisado (sin marco).
\newcommand{\figuraApaisadaSinMarco}[5]{%
  \begin{figure}%
    \begin{center}%
    \includegraphics[angle=90,height=#1\textheight,#5]{#2}%
    \caption{#3}%
    \label{#4}%
    \end{center}%
  \end{figure}%
}
% Para las tablas
\newcommand{\otoprule}{\midrule [\heavyrulewidth]}
%
% Nuevo comando para tablas pequeñas (menos de una página).
\newcommand{\tablaSmall}[5]{%
 \begin{table}[H]
  \begin{center}
   \rowcolors {2}{gray!35}{}
   \begin{tabular}{#2}
    \toprule
    #4
    \otoprule
    #5
    \bottomrule
   \end{tabular}
   \caption{#1}
   \label{tabla:#3}
  \end{center}
 \end{table}
}

%
% Nuevo comando para tablas pequeñas (menos de una página).
\newcommand{\tablaSmallSinColores}[5]{%
 \begin{table}[H]
  \begin{center}
   \begin{tabular}{#2}
    \toprule
    #4
    \otoprule
    #5
    \bottomrule
   \end{tabular}
   \caption{#1}
   \label{tabla:#3}
  \end{center}
 \end{table}
}

\newcommand{\tablaApaisadaSmall}[5]{%
\begin{landscape}
  \begin{table}
   \begin{center}
    \rowcolors {2}{gray!35}{}
    \begin{tabular}{#2}
     \toprule
     #4
     \otoprule
     #5
     \bottomrule
    \end{tabular}
    \caption{#1}
    \label{tabla:#3}
   \end{center}
  \end{table}
\end{landscape}
}

%
% Nuevo comando para tablas grandes con cabecera y filas alternas coloreadas en gris.
\newcommand{\tabla}[6]{%
  \begin{center}
    \tablefirsthead{
      \toprule
      #5
      \otoprule
    }
    \tablehead{
      \multicolumn{#3}{l}{\small\sl continúa desde la página anterior}\\
      \toprule
      #5
      \otoprule
    }
    \tabletail{
      \hline
      \multicolumn{#3}{r}{\small\sl continúa en la página siguiente}\\
    }
    \tablelasttail{
      \hline
    }
    \bottomcaption{#1}
    \rowcolors {2}{gray!35}{}
    \begin{xtabular}{#2}
      #6
      \bottomrule
    \end{xtabular}
    \label{tabla:#4}
  \end{center}
}

%
% Nuevo comando para tablas grandes con cabecera.
\newcommand{\tablaSinColores}[6]{%
  \begin{center}
    \tablefirsthead{
      \toprule
      #5
      \otoprule
    }
    \tablehead{
      \multicolumn{#3}{l}{\small\sl continúa desde la página anterior}\\
      \toprule
      #5
      \otoprule
    }
    \tabletail{
      \hline
      \multicolumn{#3}{r}{\small\sl continúa en la página siguiente}\\
    }
    \tablelasttail{
      \hline
    }
    \bottomcaption{#1}
    \begin{xtabular}{#2}
      #6
      \bottomrule
    \end{xtabular}
    \label{tabla:#4}
  \end{center}
}

%
% Nuevo comando para tablas grandes sin cabecera.
\newcommand{\tablaSinCabecera}[5]{%
  \begin{center}
    \tablefirsthead{
      \toprule
    }
    \tablehead{
      \multicolumn{#3}{l}{\small\sl continúa desde la página anterior}\\
      \hline
    }
    \tabletail{
      \hline
      \multicolumn{#3}{r}{\small\sl continúa en la página siguiente}\\
    }
    \tablelasttail{
      \hline
    }
    \bottomcaption{#1}
  \begin{xtabular}{#2}
    #5
   \bottomrule
  \end{xtabular}
  \label{tabla:#4}
  \end{center}
}



\definecolor{cgoLight}{HTML}{EEEEEE}
\definecolor{cgoExtralight}{HTML}{FFFFFF}

%
% Nuevo comando para tablas grandes sin cabecera.
\newcommand{\tablaSinCabeceraConBandas}[5]{%
  \begin{center}
    \tablefirsthead{
      \toprule
    }
    \tablehead{
      \multicolumn{#3}{l}{\small\sl continúa desde la página anterior}\\
      \hline
    }
    \tabletail{
      \hline
      \multicolumn{#3}{r}{\small\sl continúa en la página siguiente}\\
    }
    \tablelasttail{
      \hline
    }
    \bottomcaption{#1}
    \rowcolors[]{1}{cgoExtralight}{cgoLight}

  \begin{xtabular}{#2}
    #5
   \bottomrule
  \end{xtabular}
  \label{tabla:#4}
  \end{center}
}




\graphicspath{ {./img/} }

% Capítulos
\chapterstyle{bianchi}
\newcommand{\capitulo}[2]{
	\setcounter{chapter}{#1}
	\setcounter{section}{0}
	\chapter*{#2}
	\addcontentsline{toc}{chapter}{#2}
	\markboth{#2}{#2}
}

% Apéndices
\renewcommand{\appendixname}{Apéndice}
\renewcommand*\cftappendixname{\appendixname}

\newcommand{\apendice}[1]{
	%\renewcommand{\thechapter}{A}
	\chapter{#1}
}

\renewcommand*\cftappendixname{\appendixname\ }

% Formato de portada
\makeatletter
\usepackage{xcolor}
\newcommand{\tutor}[1]{\def\@tutor{#1}}
\newcommand{\course}[1]{\def\@course{#1}}
\definecolor{cpardoBox}{HTML}{E6E6FF}
\def\maketitle{
  \null
  \thispagestyle{empty}
  % Cabecera ----------------
\noindent\includegraphics[width=\textwidth]{cabecera}\vspace{1cm}%
  \vfill
  % Título proyecto y escudo informática ----------------
  \colorbox{cpardoBox}{%
    \begin{minipage}{.8\textwidth}
      \vspace{.5cm}\Large
      \begin{center}
      \textbf{TFG del Grado en Ingeniería Informática}\vspace{.6cm}\\
      \textbf{\LARGE\@title{}}
      \end{center}
      \vspace{.2cm}
    \end{minipage}

  }%
  \hfill\begin{minipage}{.20\textwidth}
    \includegraphics[width=\textwidth]{escudoInfor}
  \end{minipage}
  \vfill
  % Datos de alumno, curso y tutores ------------------
  \begin{center}%
  {%
    \noindent\LARGE
    Presentado por \@author{}\\ 
    en Universidad de Burgos --- \@date{}\\
    Tutor: \@tutor{}\\
  }%
  \end{center}%
  \null
  \cleardoublepage
  }
\makeatother


% Datos de portada
\title{Aplicación para el Acceso a Plataformas de ELearning desde Dispositivos Móviles \\Documentación Técnica}
\author{Daniel Puente Gabarri}
\tutor{María Belén Vaquerizo García y Bruno Baruque Zanón}
\date{\today}

\begin{document}

\maketitle



\cleardoublepage



%%%%%%%%%%%%%%%%%%%%%%%%%%%%%%%%%%%%%%%%%%%%%%%%%%%%%%%%%%%%%%%%%%%%%%%%%%%%%%%%%%%%%%%%



\frontmatter


\clearpage

% Indices
\tableofcontents

\clearpage

\listoffigures

\clearpage

\listoftables

\clearpage

\mainmatter

\appendix

\apendice{Plan de Proyecto Software}

\section{Introducción}

A continuación, en esta sección se encuentra toda aquella información relacionada con la gestión del proyecto. Esta información no solamente va a estar relacionada con el desarrollo del producto pedido, sino que englobará otros aspectos relevantes como la definición de los riesgos o costes del plan proyecto.
Como se ha comentado anteriormente en la sección 4.3.3, la metodología de trabajo utilizada para este proyecto es SCRUM. Para llevar cabo al desarrollo del proyecto se ha divido el desarrollo del mismo en Sprints o Iteraciones de dos semanas. Cabe destacar que algún Sprint varia en esta duración establecida debido a la disponibilidad del equipo para poder llevar a cabo la reunión al finalizar cada iteración.


\section{Planificación temporal}

Como se comentaba, se ha decidido dividir el desarrollo del proyecto en Sprints de dos semanas de duración. Gracias a esta duración se podrá compaginar de una manera más eficiente el desarrollo del proyecto de final de grado junto con el resto de asignaturas de este cuatrimestre. En este cuatrimestre el número de asignaturas cursadas son tres junto con el trabajo de fin de grado. Además, se ha estipulado que el número de horas de dedicación por iteración sea de alrededor de cuarenta, este es el principal motivo por el cual se ha decidido que la duración de cada iteración sea de quince días para poder compaginar el desarrollo del proyecto junto con el resto de asignaturas.

\subsection{Sprint 1. Inicio del proyecto (7 Feb 2017 a 16 Feb 2017)}

Este día tuvo lugar la primera reunión para la explicación sobre el producto a desarrollar. Una vez explicado el mismo, se decidió aceptar este proyecto y se comenzó con la definición de los principales pilares a superar para construir el producto. Una vez definidas estas bases se decidió realizar en este Sprint las siguientes tareas.

\begin{itemize}

	\item \textbf{Cargar código fuente:} este primer commit fue realizado por los tutores para incluir el proyecto de partida.
	
	\item \textbf{Objetivos del proyecto:} esta tarea ha sido realizada por el alumno sobre la documentación a entregar. En dicho apartado del documento se definirán las bases principales del proyecto junto con sus objetivos más destacados.
	\item \textbf{Instalar proyecto de partida:} en tarea fue realizada por el alumno, para ello se ha realizado una lectura de la documentación del proyecto de partida y su posterior instalación para conocer al máximo la lógica a implementar.
	
	\item \textbf{Generar el Product Backlog:}  esta tarea fue realizada por el alumno en la memoria. En esta sección se llevó a cabo la definición de los principales requisitos funcionales las cuales derivarán en las historias de usuario a realizar en los diferentes Sprint.
	
	\item \textbf{Aprendizaje PHP:} esta tarea fue la más costosa de esta iteración ya que se desconocía cualquier conocimiento acorde a este lenguaje. Para ello se realizado una lectura de la documentación proporcionada en la propia web de PHP \cite{wiki:phpdoc} junto con otras fuentes online como tutoriales sobre PHP desde cero.

\end{itemize}


A continuación, en la siguiente Ilustración se incluye un gráfico Burndown resumiendo el Sprint, cabe destacar que en esta ocasión solamente se aprecia que todas las tareas fueron finalizadas el mismo día, aunque esto no fue así. Este error nace fruto de un desconocimiento previo de la herramienta, al no incluir en la configuración de la creación de dicha gráfico las tareas que se encuentren dentro del tablero en la sección Done. Posteriormente, una vez localizado y subsanado este error en las siguientes iteraciones este fallo no se produce.

\imagen{Sprint1}{Gráfico Burndown Sprint 1 - Inicio del proyecto.}


\subsection{Sprint 2. Estudio de refactorización del servidor (16 Feb 2017 a 2 Mar 2017)}

Al comienzo de esta iteración se realizó la revisión de las tareas a realizar en el Sprint anterior y una vez revisadas se dio por finalizado el Sprint 1. 
En esta iteración se decidió realizar las siguientes tareas:

\begin{itemize}
	
	\item Estudio de inclusión de un Framework: esta tarea fue realizada por el alumno y se realizado un estudio de diferentes alternativas como: CakePHP \cite{wiki:cakephp}, Laravel \cite{wiki:laravel}, Symphony \cite{wiki:symphony}, Slim \cite{wiki:slim}. No obstante, la decisión ha sido no utilizar ningún Framework al tratarse de un proyecto pequeño en dimensiones y al añadir mayor complejidad a la resolución del mismo.
	
	\item •	Listar los métodos a implementar en el API: esta tarea ha sido realizada por el alumno y en ella se listan las principales lógicas a implementar sobre la base de datos.
	
	\item •	Estudio de diferencias entre REST y SOAP: esta tarea también fue realizada por el alumno, para ello se realizó una correcta lectura y posterior documentación en el apartado Soap y Rest. Fruto de este estudio se ha decidido utilizar REST como arquitectura para facilitar el uso y tratamiento de las peticiones desde el programa.
	
	\item •	Documentación: Aspectos Teóricos / Técnicas y Herramientas: esta tarea consiste en llevar a cabo una documentación en la memoria sobre estos aspectos en la sección Conceptos teóricos y Técnicas y herramientas.


\end{itemize}

En la siguiente Ilustración podemos observar el gráfico Burndown correspondiente a esta iteración.

\imagen{Sprint2}{Gráfico Burndown Sprint 2 - Estudio de refactorización del servidor. }

\subsection{Sprint 3. Diseño del BackEnd (2 Mar 2017 a 15 Mar 2017)}

Al comienzo de esta iteración se realizó la revisión de las tareas a realizar en el Sprint anterior y una vez revisadas se dio por finalizado el Sprint 2. 
En esta iteración se decidió realizar las siguientes tareas:

\begin{itemize}

	\item \textbf{Implementación de los métodos del API:} esta tarea consiste en realizar la implementación del APIREST, para poder comunicarnos con el controlador del proyecto de partida.

	\item \textbf{Diseño arquitectura del BackEnd:} esta tarea consiste en construir los diagramas sobre el BackEnd.

	\item \textbf{Anexo 1 – Planificación del proyecto:} esta tarea consiste en añadir a la documentación este anexo.
\end{itemize}

En la siguiente Ilustración podemos observar el gráfico Burndown correspondiente a esta iteración.

\imagen{Sprint3}{Gráfico Burndown Sprint 3 - Diseño del BackEnd. }

\subsection{Sprint 4. Interconexión FrontEnd-BackEnd (15 Mar 2017 a 29 Mar 2017)}

Al comienzo de esta iteración se realizó la revisión de las tareas a realizar en el Sprint anterior y una vez revisadas se dio por finalizado el Sprint 3. 
En esta iteración se decidió realizar las siguientes tareas:

\begin{itemize}

	\item \textbf{Test unitarios sobre el BackEnd}: en esta tarea se realizo una documentación sobre las pruebas unitarias en PHP. Posteriormente, se realizaron pruebas unitarias sobre las nuevas clases creadas en el APIREST.
	\item \textbf{Leer documentación LTI}: en esta tarea se realizo una lectura sobre el funcionamiento de LTI para su posterior utilización en la aplicación
	\item \textbf{Prueba conexión Android - Moodle}: esta tarea fue la mas complicada del Sprint y posiblemente una de las mas costosas del proyecto. Después de mucha lectura se decidió utilizar los Web service de Moodle para obtener la información de la plataforma desde Android.
	\item \textbf{Documentacion, Anexo Diseño}: en esta tarea se incluyo en la memoria los diagramas de clases y paquetes del APIREST.
	\item \textbf{Definir ProductBackLog}: en esta tarea se definió el conjunto de historias de usuario a desarrollar.

\end{itemize}

En la siguiente Ilustración podemos observar el gráfico Burndown correspondiente a esta iteración.

\subsection{Sprint 5. Recuperación de datos desde el BackEnd (29 Mar 2017 a 20 Abr 2017)}

Al comienzo de esta iteración se realizó la revisión de las tareas a realizar en el Sprint anterior y una vez revisadas se dio por finalizado el Sprint 4. 
En esta iteración se decidió realizar las siguientes tareas:

\begin{itemize}

	\item \textbf{Implementar conexión LTI desde Android}: en esta tarea se realizó la obtención del token del usuario que ha iniciado sesión en la aplicación.
	\item \textbf{MockUp del la App Android}: en esta tarea se realizo una construcción genérica de la app.


\end{itemize}

En la siguiente Ilustración podemos observar el gráfico Burndown correspondiente a esta iteración.

\subsection{Sprint 6. Recuperación de Datos desde BackEnd (II) (20 Abr 2017 a 4 May 2017)}

Al comienzo de esta iteración se realizó la revisión de las tareas a realizar en el Sprint anterior y una vez revisadas se dio por finalizado el Sprint 5. 
En esta iteración se decidió realizar las siguientes tareas:

\begin{itemize}

	\item \textbf{Diseño de la arquitectura del FrontEnd}: en esta tarea se incluyó en la memoria el diagrama de clases y paquetes de la aplicación Android.
	\item \textbf{Obtención datos, herramienta externa}: en esta tarea se realizó un estudio sobre los diferentes métodos del Web service de Moodle necesarios para la obtención de los cuestionarios de QuickTest.
	\item \textbf{HU1 - Iniciar sesión}:en esta tarea se proporcionó un diseño al login más cercano al producto final.
	\item \textbf{HU2 - Cuestionarios a resolver}: en esta tarea se obtuvieron los cuestionarios de QuickTest asociados al usuario logeado.
	\item \textbf{Añadir aspectos relevantes hasta el momento en el manual del programador}: en esta tarea se completo esta sección en la memoria.
	\item \textbf{Casos de Uso}: en esta tarea se completo esta sección en la memoria.
	\item \textbf{Obtener datos desde QuickTest}: en esta se incluyó en el APIREST la funcionalidad de obtener un cuestionario por identificador, ya que esta funcionalidad no se encontraba desarrollada en el controlador,
	\item \textbf{Incluir nuevas bibliotecas en Manual Programador}: en esta tarea se añadieron una explicación sobre las librerías Gson y Volley.


\end{itemize}

En la siguiente Ilustración podemos observar el gráfico Burndown correspondiente a esta iteración.

\subsection{Sprint 7. Grabado de datos desde el BackEnd (4 May 2017 a 18 May 2017)}

Al comienzo de esta iteración se realizó la revisión de las tareas a realizar en el Sprint anterior y una vez revisadas se dio por finalizado el Sprint 6. 
En esta iteración se decidió realizar las siguientes tareas:

\begin{itemize}

	\item \textbf{HU3 - Resolver un cuestionario}: en esta tarea realizó el desarrollo de esta historia de usuario. Destacar que solamente se permitió que el alumno resolviera un cuestionario sin la posibilidad de los comodines.
	\item \textbf{HU4 - Finalizar un cuestionario}: en esta tarea se permitió que el alumno pueda finalizar un cuestionario desde la aplicación.
	\item \textbf{HU7 - Podrá cerrar sesión}: en esta tarea se realizó el desarrollo de esta historia de usuario en la aplicación.
	\item \textbf{Completar Aspectos Relevantes (Memoria)}: en esta tarea se completo esta sección en la memoria.
	\item \textbf{Envío de respuestas desde la app a Quicktest}: en esta tarea se permitió que al finalizar un cuestionario, estos datos sean enviados al APIREST para su almacenamiento.


\end{itemize}

En la siguiente Ilustración podemos observar el gráfico Burndown correspondiente a esta iteración.


\subsection{Sprint 8. Ampliación del Front End. (18 May 2017 a 8 Jun 2017)}

Al comienzo de esta iteración se realizó la revisión de las tareas a realizar en el Sprint anterior y una vez revisadas se dio por finalizado el Sprint 7. 
En esta iteración se decidió realizar las siguientes tareas:

\begin{itemize}

	\item \textbf{HU5 - Calificación obtenida}: en esta tarea se incluyó en la aplicación una sección donde el usuario pudiera ver la calificación obtenida en un cuestionario ya resuelto.
	\item \textbf{HU6 - Revisar un cuestionario}: en esta tarea se incluyó en la aplicación una sección donde el usuario pudiera ver una tabla con información sobre el cuestionario.
	\item \textbf{Implementación de Funcionalidad de Comodines}: en esta tarea se completó la HU3 al permitir al alumno utilizar comodines.


\end{itemize}

En la siguiente Ilustración podemos observar el gráfico Burndown correspondiente a esta iteración.

\imagen{Sprint8}{Gráfico Burndown Sprint 8 - Ampliación del Front End. }

\subsection{Sprint 9. Finalizar FrontEnd y test. (8 Jun 2017 a 19 Jun 2017)}

Al comienzo de esta iteración se realizó la revisión de las tareas a realizar en el Sprint anterior y una vez revisadas se dio por finalizado el Sprint 8. 
En esta iteración se decidió realizar las siguientes tareas:

\begin{itemize}

	\item 


\end{itemize}

En la siguiente Ilustración podemos observar el gráfico Burndown correspondiente a esta iteración.

\imagen{Sprint8}{Gráfico Burndown Sprint 8 - Ampliación del Front End. }



\section{Estudio de viabilidad}

Uno de los objetivos principales de cualquier proyecto por no clasificarlo como el más importante es conocer si este proyecto es viable. Para ello se debe realizar un estudio sobre la viabilidad del mismo. Es por esto que se deberá realizar este estudio sobre dos enfoques: \emph{económico y legal}. Y gracias a estos sabremos si el proyecto es rentable y viable.

\subsection{Viabilidad económica}

Relacionado con lo anteriormente expuesto en esta sección se va realizar un estudio sobre la viabilidad del proyecto en el enfoque económico. Para ello se va realizar un análisis sobre los diferentes costes que entran en juego y son necesarios para el desarrollo del mismo junto con el mantenimiento una vez desarrollado.

\subsection{Coste Hardware}

Para el desarrollo del proyecto y posterior corrección de ciertos problemas que puedan surgir será necesario:

\begin{itemize}

    \item \textbf{Un ordenador potente en prestaciones}, ya que el IDE de Android Studio es bastante pesado en tareas principalmente de emulación. 
	
	\item \textbf{Un dispositivo móvil}, es decir un Smartphone con S.O Android.

\end{itemize}

En este caso se utilizarán las herramientas propias del alumno encargado de realizar el proyecto, que son las siguientes:

\begin{itemize}

	\item \textbf{Asus GL552VW-DM142T:} coste actual 1199 \euro sin S.O.
	
	\item \textbf{OnePlus 2:} coste actual 319 \euro.

\end{itemize}


Ambos componentes son bastante potentes y por lo tanto no ocasionarán ninguna problemática de que el hardware se quede anticuado en prestaciones en posteriores años. Para redondear los posteriores cálculos estimaremos que el tiempo de vida de ambos componentes será de 5 años, es decir 60 meses. Aunque seguramente su tiempo de vida podría alargarse aún más. Además, al realizarse el proyecto en un cuatrimestre la duración de la utilización de dichos componentes será de 4 meses.

\tablaSmall{Coste Hardware}{l l}{herramientasportipodeuso}
{ 
Tiempo de amortización & 60 meses\\
Coste de los componentes & 1199,00 \euro + 319,00 \euro = 1518,00 \euro \\
Coste de amortización/mes & 1518,00 \euro / 60 meses = 25,3 \euro /mes \\
Coste final & 25,3 \euro /mes * 4 meses = 101,2 \euro \\

}





\subsection{Coste Software}

Para la realización del proyecto y su posterior mantenimiento será necesario lo siguiente:

\tablaSmall{Coste Software 1}{l l l }{coste-1}
{ \multicolumn{1}{l}{Software} & Licencia & Coste \\}{ 
Microsoft Windows 10 & Home & 135,00 \euro \\
Android Studio & Licencia Apache 2.0 \cite{wiki:astudio} & 0,00 \euro \\
XAMPP & GNU \cite{wiki:xampp} & 0,00 \euro \\
GitHub & X & X \\	
TortoiseSVN & GNU General Public License \cite{wiki:tortoisesvn} & 0,00 \euro \\
Moodle & GNU GPL \cite{wiki:moodle} & 0,00 \euro \\
PhpStorm & Estudiante & 0,00 \euro \\
SonarQube & LGPL & 0,00 \euro \\
Advanced REST Client & X & X \\
Microsoft Office & Hogar & 99,00 \euro \\
& & 234,00 \euro \\

}

En este caso la licencia de Microsoft Office es para un único año de duración, pero la licencia del S.O no tiene caducidad por lo que solamente tendremos en cuenta una duración anual.


\tablaSmall{Coste Software 2}{c c}{H}
{ 
Tiempo de amortización & 12 meses\\
Coste de los componentes & 135,00 \euro + 99,00 \euro = 234,00 \euro \\
Coste de amortización/mes & 234,00 \euro / 12 meses = 19,5 \euro /mes \\
Coste final & 19,5 \euro /mes * 4 meses = 78,00 \euro \\
}


\subsection{Coste de instalación}

Este coste será nulo y no se tendrá en cuenta ya que la instalación de los distintos componentes será realizada por el propio alumno.

\subsection{Coste de aprendizaje}

Este proyecto será desarrollado por el alumno por lo que este coste de realización será un sueldo mensual a dicho alumno por llevar a cabo el diseño, implementación y testeo del mismo. Como se comentaba anteriormente se ha estipulado que el número de horas dedicadas al desarrollo del mismo será de 40 horas cada dos semanas.

\tablaSmall{Coste de aprendizaje}{c c}{CostesDeAprendizaje}
{ 
Duración & 4 meses * 4 semanas = 16 semanas\\
Salario/hora & 10,00 \euro /hora\\
Nºhoras/semana & 20 horas/semana \\
Coste final & 10,00 \euro /hora * 20 horas/semana * 16 semanas = 3.200,00 \euro \\
}

\subsection{Coste de soporte}

Este coste sería realizado por un hipotético técnico que se ocuparía de resolver cualquier incidencia en la utilización del producto desarrollado. Al ser una aplicación que, si resulta viable, será publicada en el Play Store no existirá ningún técnico encargado de subsanar directamente los errores a los usuarios, sino que serán los propios operarios de dicha página los que se remitirán cualquier incidencia.
Es por esto que coste de soporte será nulo y por tanto no se tendrá en cuenta.

\subsection{Coste de mantenimiento}

En este caso será el propio alumno el encargado de realizar el mantenimiento de la aplicación por lo que esta tarea se encuentra remunerada dentro de su salario. No obstante, cabe destacar que la distribución desde esta App no incluye costes de almacenamiento de la información resultantes, es decir, las propias instituciones que compren el producto tendrán que realizar un mantenimiento propio de sus servidores.

\subsection{Otros costes}

Además, también hay que tener en cuenta otros costes como los derivados del \emph{material mobiliario, de oficina, documentación, electricidad, Internet}.

\tablaSmall{Otros costes}{l c c }{Otros Costes}
{ \multicolumn{1}{l}{Tipos de coste} & Coste & Total \\}{ 
Infraestuctura & 10,00 \euro /mes &  10,00 \euro /mes * 4 meses = 40,00 \euro \\
Documentación & 36,10 \euro \cite{wiki:book}  & 36,10 \euro \\
Electricidad & 35,00 \euro & 35 \euro /mes * 4 = 140 \euro \\
Internet & X & 37,90 \euro \euro /mes * 4 = 151,60 \euro \\	

}

\subsection{Total}


\tablaSmall{Costes totales}{l c}{H}
{ \multicolumn{1}{l}{Tipos de coste} & Total \\}{ 

Coste Hardware & 101,2 \euro \\
Coste Software & 78,00 \euro \\
Coste de instalación & 0,00 \euro \\
Coste de aprendizaje & 3.200,00 \euro \\
Coste de soporte & 0,00 \euro \\
Coste de mantenimiento & 367,70 \euro \\
Otros coste & 367,70 \euro \\
& 3.746,90 \euro \\

}

\subsection{Análisis coste-beneficio}

Como ya se ha explicado, este proyecto nace de la base fundamental de otro proyecto. Este proyecto fue desarrollado con una base y un fin muy bien definido y fijado, la herramienta QuickTest es compatible con cualquier plataforma que cumpla el estándar LTI. Fruto de esta idea se desarrolló este proyecto como un complemento de resolución de cuestionarios para Moodle. Es por esto que este análisis de coste-beneficio se va realizar para Moodle.
Moodle cuenta actualmente con más de 100 millones de usuarios en todo el mundo, lo que la convierte en la plataforma digital más utilizada en todo mundo por la comunidad educativa \cite{wiki:mest} \cite{wiki:moodlestas}. Además, se encuentra registrada en 234 países, en la que se encuentra España en el segundo puesto con 7.152 inscripciones.
Como el principal objetivo de esta aplicación es la utilización de la misma para un entorno universitario vamos a reducir estos cálculos de inscripciones al número de universidades en España. Actualmente, el número de universidades tanto públicas como privadas es de 82. 
Al ser el proyecto compatible con Moodle y al ser este, la plataforma digital más influyente en este sector, cabe destacar que no deberíamos de tener ningún inconveniente en vender licencias de este software en el PlayStore. Sería una suposición muy idealizada que todas estas comprarán nuestro producto por lo que realizaremos los cálculos estimando que solamente 40 universidades deciden adquirir la licencia.
Ante este estudio preliminar, se decide vender cada licencia por 100,00 \euro. Este precio solamente incluirá el producto junto con la documentación necesarios para su correcta utilización, es decir, no se incluye un soporte de almacenamiento de la información generada por su utilización. Además, esta distribución de licencia se hará con carácter comercial \cite{wiki:liccom}. 
Finalmente, cabe destacar que valoraremos este análisis como factible ya que únicamente distribuyendo estas 40 licencias recuperaríamos la inversión realizada.



\subsection{Viabilidad legal}

Al utilizar herramientas que se encuentran bajo licencias de software libre la distribución de esta herramienta no ocasionará ninguna problemática legal.
No obstante, la herramienta fruto de este proyecto se distribuirá con la licencia Android License Verification Library \cite{wiki:androidlic} propia de PlayStore y utilizada por cualquier aplicación que resida en la misma para evitar posibles pirateos.
Como ya se comentaba en la sección anterior este producto se distribuirá como licencia comercial, ya que todos los activos utilizados para el desarrollo del mismo poseen licencia GNU.




\apendice{Especificación de Requisitos}

\section{Introducción}

Como ya se ha comentado anteriormente, se va utilizar la metodología ágil SCRUM. Es por esto que no vamos a tratar las tareas a realizar como requisitos, sino como historias de usuario. Todas estas historias de usuarios formarán parte del Product Backlog y serán asignadas a determinadas iteraciones durante el desarrollo del proyecto.
A continuación, se incluye un listado sobre las historias de usuarios y sus correspondientes diagramas de casos de uso.


\section{Product Backlog}

En esta sección se incluyen todas aquellas historias de usuarios necesarias para los alumnos y los profesores.

\subsection{Historias de usuario de ambos roles}

\subsubsection{Historia de usuario 1: Iniciar sesión}\label{subsec:HU1}

\tablaSmall{Historia de usuario 1}{l l}{coste-1}
{ \multicolumn{1}{l}{HU1}\\}{ 
Título & Como un usuario autenticado, podrá iniciar sesión. \\
Rol & Alumno o profesor autenticado. \\
Descripción & El usuario introducirá sus credenciales y se le mostrará \\\
& aquellos cursos en los que se encuentre matriculado. \\
Precondiciones & - El usuario debe estar previamente registrado. \\\
			   & - Haber iniciado sesión. \\
}

\subsubsection{Historia de usuario 2: Cerrar sesión}\label{subsec:HU2}

\tablaSmall{Historia de usuario 2}{l l}{coste-1}
{ \multicolumn{1}{l}{HU2}\\}{ 
Título & Como un usuario autenticado, podrá cerrar sesión.\\
Rol & Alumno o profesor autenticado. \\
Descripción & Siempre y cuando el usuario se encuentre logeado, podrá\\\
& cerrar sesión.\\
Precondiciones & - El usuario debe estar previamente registrado.\\\
& - Haber iniciado sesión. \\
}

\subsubsection{Historia de usuario 3: Recordar campos}\label{subsec:HU3}

\tablaSmall{Historia de usuario 3}{l l}{coste-1}
{ \multicolumn{1}{l}{HU3}\\}{ 
Título & Como un usuario autenticado, podrá recordar los campos.\\
Rol & Alumno o profesor autenticado. \\
Descripción & El sistema permitirá recordar los campos con los \\\
& que el usuario inicia sesión.\\
Precondiciones & - El usuario debe estar previamente registrado.\\\
& - Haber iniciado sesión. \\
}

\subsubsection{Historia de usuario 4: Cursos}\label{subsec:HU4}

\tablaSmall{Historia de usuario 4}{l l}{coste-1}
{ \multicolumn{1}{l}{HU4}\\}{ 
Título & Como un usuario autenticado, podrá ver sus cursos.\\
Rol & Alumno o profesor autenticado. \\
Descripción & El sistema mostrará los cursos en los que se encuentra \\\ & matriculado.\\
Precondiciones & - El usuario debe estar previamente registrado.\\\
& - Haber iniciado sesión. \\
}

\subsection{Historias de usuario de los alumnos}

\subsubsection{Historia de usuario 5: Cuestionarios}\label{subsec:HU5}

\tablaSmall{Historia de usuario 5}{l l}{coste-1}
{ \multicolumn{1}{l}{HU5}\\}{ 
Título & Como un usuario autenticado, el sistema notificará sobre\\\
& los cuestionarios a resolver. \\
Rol & Alumno autenticado. \\
Descripción & Una vez el usuario haya iniciado sesión obtendrá todos\\\  & aquellos cuestionarios en los que se encuentre asignado. \\
Precondiciones & - El usuario debe estar previamente registrado.\\\
& - Haber iniciado sesión. \\\
& - Existan cuestionarios en aquellos cursos en los que\\\
& se encuentra matriculado el alumno. \\
}

\subsubsection{Historia de usuario 6: Resolver cuestionario}\label{subsec:HU6}

\tablaSmall{Historia de usuario 6}{l l}{coste-1}
{ \multicolumn{1}{l}{HU6}\\}{ 
Título & Como un usuario autenticado, podrá resolver un \\\
& cuestionario. \\
Rol & Alumno autenticado. \\
Descripción & Una vez el usuario haya iniciado sesión podrá resolver \\\ & aquellos cuestionarios en los que se encuentren asignado. \\
Precondiciones & - El usuario debe estar previamente registrado.\\\
& - Haber iniciado sesión. \\\
& - Existan cuestionarios en aquellos cursos en los que se\\\
& encuentra matriculado el alumno. \\
}

\subsubsection{Historia de usuario 7: Finalizar cuestionario}\label{subsec:HU7}

\tablaSmall{Historia de usuario 7}{l l}{coste-1}
{ \multicolumn{1}{l}{HU7}\\}{ 
Título & Como un usuario autenticado, podrá finalizar un \\\ & cuestionario. \\
Rol & Alumno autenticado. \\
Descripción & El usuario podrá finalizar aquel cuestionario que\\\
& este resolviendo. \\
Precondiciones & - El usuario debe estar previamente registrado.\\\
& - Haber iniciado sesión. \\\
& - Existan cuestionarios en aquellos cursos en los que se \\\
& encuentra matriculado el alumno. \\\
& - Haber iniciado la resolución de un cuestionario. \\
}

\subsubsection{Historia de usuario 8: Calificación obtenida}\label{subsec:HU8}

\tablaSmall{Historia de usuario 8}{l l}{coste-1}
{ \multicolumn{1}{l}{HU8}\\}{ 
Título & Como un usuario autenticado, el sistema notificará \\\
& sobre la calificación obtenida. \\
Rol & Alumno autenticado. \\
Descripción & Una vez el usuario finalice el cuestionario, el sistema \\\
& proporcionará una retroalimentación sobre el cuestionario \\\
& resuelto.\\
Precondiciones & - El usuario debe estar previamente registrado.\\\
& - Haber iniciado sesión. \\\
& - Existan cuestionarios en aquellos cursos en los que se\\\
&  encuentra  matriculado el alumno. \\\
& - Haber iniciado la resolución de un cuestionario. \\\
& - Haber finalizado un cuestionario. \\
}

\subsubsection{Historia de usuario 9: Revisar cuestionario resuelto}\label{subsec:HU9}

\tablaSmall{Historia de usuario 9}{l l}{coste-1}
{ \multicolumn{1}{l}{HU9}\\}{ 
Título & Como un usuario autenticado, podrá revisar un \\\
& cuestionario resuelto.\\\
Rol & Alumno autenticado. \\
Descripción & El usuario podrá revisar la retroalimentación de un  \\\ & cuestionario que ya esté finalizado. \\
Precondiciones & - El usuario debe estar previamente registrado.\\\
& - Haber iniciado sesión. \\\
& - Existan cuestionarios en aquellos cursos en los que se\\\
&  encuentra matriculado el alumno. \\\
& - Haber iniciado la resolución de un cuestionario. \\\
& - Haber finalizado un cuestionario. \\
}

\subsection{Historias de usuario de los profesores}

\subsubsection{Historia de usuario 10: Cuestionarios}\label{subsec:HU10}

\tablaSmall{Historia de usuario 10}{l l}{coste-1}
{ \multicolumn{1}{l}{HU10}\\}{ 
Título & Como un usuario autenticado, podrá ver los cuestionarios \\\ & a los que se enfrentan sus alumnos.\\
Rol & Profesor autenticado. \\
Descripción & El sistema mostrará al profesor los cuestionarios a los \\\
& que se enfrentan sus alumnos.\\
Precondiciones & - El usuario debe estar previamente registrado.\\\
& - Haber iniciado sesión. \\
}

\subsubsection{Historia de usuario 11: Calificaciones}\label{subsec:HU11}

\tablaSmall{Historia de usuario 11}{l l}{coste-1}
{ \multicolumn{1}{l}{HU11}\\}{ 
Título & Como un usuario autenticado, podrá ver las calificaciones \\\ & de los alumnos en un cuestionario.\\
Rol & Profesor autenticado. \\
Descripción & El sistema mostrará al profesor la calificación de cada \\\ & alumno en cada cuestionario.\\
Precondiciones & - El usuario debe estar previamente registrado.\\\
& - Haber iniciado sesión. \\
}

\subsubsection{Historia de usuario 12: Ver cuestionarios}\label{subsec:HU12}

\tablaSmall{Historia de usuario 12}{l l}{coste-1}
{ \multicolumn{1}{l}{HU12}\\}{ 
Título & Como un usuario autenticado, podrá ver los cuestionarios.\\
Rol & Profesor autenticado. \\
Descripción & El sistema mostrará al profesor los cuestionarios a los que se \\\ & enfrentan sus alumnos.\\
Precondiciones & - El usuario debe estar previamente registrado.\\\
& - Haber iniciado sesión. \\
}

\section{Diagrama de casos de uso}

En la siguiente ilustración podemos ver el diagrama de los casos de uso de nuestro sistema:

\imagen{DiagramaCasosDeUso}{Diagrama de casos de uso del alumno.}

\imagen{DiagramaCasosDeUsoProfesor}{Diagrama de casos de uso del profesor.}

\subsection{Caso de uso 1: Iniciar sesión}

\tablaSmall{Caso de uso 1: Iniciar sesión}{l l}{coste-1}
{ \multicolumn{1}{l}{CU1}\\}{ 
Título & Iniciar sesión.\\
Descripción & El sistema permitirá al usuario autenticado iniciar sesión para \\\ & acceder al sistema.\\
Secuencia & 1. El usuario introducirá los campos de nombre y contraseña. \\\
& 2. El sistema comprobará si los campos son correctos: \\\
& \hspace{0.25cm} 2.1. Si son correctos, el usuario pasa a estar logeado teniendo  \\\ & \hspace{0.25cm} pleno acceso a las funcionalidades del sistema. \\\
& \hspace{0.25cm} 2.2. Si no son correctos, el sistema notificará que los campos \\\ & son incorrectos. \\
Precondiciones & - El usuario debe estar previamente registrado en Moodle. \\
Comentarios & Este caso de uso, CU1, corresponde a la HU1. \ref{subsec:HU1} \\	
}

\subsection{Caso de uso 2: Olvidar campos}

\tablaSmall{Caso de uso 2: Olvidar campos}{l l}{coste-1}
{ \multicolumn{1}{l}{CU2}\\}{ 
Título & Recordar campos.\\
Descripción & El sistema permitirá al usuario olvidar los campos con las que \\\ & el usuario inicia sesión.\\
Secuencia & 1. El usuario seleccionara en el sistema cerrar sesión. \\\
& 2. El sistema olvidará las credenciales del usuario.\\
Precondiciones & - El usuario debe estar previamente registrado en Moodle.\\\
& - El usuario debe haber iniciado sesión y haber activado la 
\\\ & opción  de recordar campos.\\
Comentarios & Este caso de uso, CU2, corresponde a la historia de usuario 2. \\\ & \ref{subsec:HU2} \\
}

\subsection{Caso de uso 3: Recordar campos}

\tablaSmall{Caso de uso 3: Recordar campos}{l l}{coste-1}
{ \multicolumn{1}{l}{CU3}\\}{ 
Título & Olvidar campos.\\
Descripción & El sistema permitirá al usuario recordar los campos con las que \\\ & el usuario inicia sesión.\\
Secuencia & 1. El sistema recordará las credenciales del usuario. \\
Precondiciones & - El usuario debe estar previamente registrado en Moodle.\\
Comentarios & Este caso de uso, CU3, corresponde a la historia de usuario 3. \\\ & \ref{subsec:HU3}\\
}

\subsection{Caso de uso 4: Mostrar cursos}

\tablaSmall{Caso de uso 4: Mostrar cuestionarios}{l l}{coste-1}
{ \multicolumn{1}{l}{CU4}\\}{ 
Título & Mostrar cursos.\\
Descripción & El sistema permitirá al usuario ver los cursos en \\\
& los que se encuentra matriculado.\\
Secuencia & 1. El sistema mostrará los cursos. \\
Precondiciones & - El usuario debe estar previamente registrado en Moodle. \\\ & - El usuario debe estar logeado en la aplicación para poder  \\\
& acceder a esta funcionalidad. \\\
& - Deben existir cuestionarios de QuickTest en cursos en los que  \\\ & el usuario se encuentre matriculado.\\
Comentarios & Este caso de uso, CU4, corresponde a la historia de usuario 4. \\\ & \ref{subsec:HU4} \\
}


\subsection{Caso de uso 5: Mostrar cuestionarios}

\tablaSmall{Caso de uso 5: Mostrar cuestionarios}{l l}{coste-1}
{ \multicolumn{1}{l}{CU5}\\}{ 
Título & Mostrar cuestionarios.\\
Descripción & El sistema permitirá al usuario ver los cuestionarios de  \\\
& un curso.\\
Secuencia & 1. El sistema mostrará los cuestionarios. \\
Precondiciones & - El usuario debe estar previamente registrado en Moodle. \\\ & - El usuario debe estar logeado en la aplicación para poder  \\\
& acceder a esta funcionalidad. \\\
& - Deben existir cuestionarios de QuickTest en cursos en los que  \\\ & el usuario se encuentre matriculado.\\
Comentarios & Este caso de uso, CU5, corresponde a la historia de usuario 5. \\\ & \ref{subsec:HU5} \\
}

\subsection{Caso de uso 6: Resolver cuestionario}

\tablaSmall{Caso de uso 6: Resolver cuestionario}{l l}{coste-1}
{ \multicolumn{1}{l}{CU6}\\}{ 
Título & Resolver cuestionario.\\
Descripción & El sistema permitirá al usuario resolver un cuestionario. \\
Secuencia & 1. El usuario seleccionará un cuestionario a resolver y comenzará \\\ & \hspace{0.25cm} a resolverlo. \\\
&  2. Una vez finalizado el cuestionario, el sistema enviará el \\\
& \hspace{0.25cm} cuestionario resuelto HU7. \ref{subsec:HU7} \\\
& 3. El sistema mostrará la retroalimentación del cuestionario \\\ & \hspace{0.25cm} resuelto HU8, HU9. \ref{subsec:HU8}\\
Precondiciones & - El usuario debe estar previamente registrado en Moodle.\\\
& - El usuario debe estar logeado en la aplicación para poder 
\\\ & acceder a esta funcionalidad. \\\
& - Deben existir cuestionarios de QuickTest en cursos en los que \\\ & el usuario se encuentre matriculado y que hayan sido resueltos.\\
Comentarios & Este caso de uso, CU6, corresponde a la historia de usuario \\\ & 6, 7, 8, 9.  \ref{subsec:HU6}\\
}

\subsection{Caso de uso 7: Mostrar cuestionarios resueltos}

\tablaSmall{Caso de uso 7: Mostrar cuestionarios resueltos}{l l}{coste-1}
{ \multicolumn{1}{l}{CU7}\\}{ 
Título & Mostrar cuestionarios resueltos.\\
Descripción & El sistema permitirá al usuario ver los cuestionarios resueltos. \\
Secuencia & 1. El sistema mostrará los cuestionarios resueltos. \\\
& \hspace{0.25cm} 1.1. Si existe algún cuestionario resuelto el usuario   \\\ & \hspace{0.25cm} podrá seleccionar un cuestionario y obtener la  \\\ & \hspace{0.25cm} retroalimentación del mismo HU8, HU9. \\
Precondiciones & - El usuario debe estar previamente registrado en Moodle.\\\
& - El usuario debe estar logeado en la aplicación para poder \\\ & acceder a esta funcionalidad. \\\
& - Deben existir cuestionarios de QuickTest en cursos en los que   \\\ & el usuario se encuentre matriculado y que hayan sido resueltos.\\
Comentarios & Este caso de uso, CU7, corresponde a la historia de usuario 6. \\\ & \ref{subsec:HU7}\\
}







\apendice{Especificación de diseño}

\section{Introducción}

A continuación, se incluye una explicación/descripción del diseño de la aplicación Android. Para ello, se dividirán los diagramas en dos enfoques: el APIREST utilizado para comunicarse con el proyecto de partida y la propia app Android que se comunicará con dicho API para la obtención y volcado de los datos.

\section{Diseño de datos}

En este apartado se incluirán los diagramas de clases y paquetes de ambos enfoques.

\subsection{Diagrama de clases}

En estos diagramas se ha decidido utilizar un criterio de color para diferenciar aquellas clases que forman parte del APIREST, en color azul y en amarillo aquellos que forman parte del proyecto de partida.

\subsubsection{APIREST}

Este paquete será el encargado de comunicarse con el Controlador del proyecto de partida. A continuación, se mostrarán de manera fraccionada cada una de las clases que lo forman y con qué clases del controlador del proyecto de partida se comunican.

\paragraph{Clase ControlAccesoProfesor}

Esta clase será la encargada de realizar el control de acceso por parte de los profesores a la aplicación web. Esta web provee a los profesores de una clave privada dentro de la aplicación para poder utilizarla en la sincronización de sus cuestionarios con Moodle. Esta clase permite registrar a un nuevo profesor o comprobar si puede iniciar sesión. Ambas lógicas se controlan mediante una petición post. Para llevar a cabo este cometido se comunican con el controlador llamando al método registrarNuevoUsuario o login respectivamente.

\begin{itemize}

	\item \textbf{Post:} esta petición será la encargada de manejar la lógica de toda la clase.
	
	\begin{list}{-}
		\item \emph{Registrarse en QuickTest:} para ello es necesario incluir dentro del propio cuerpo de la petición el email y contraseña del usuario. La petición es la siguiente: \url{http://localhost/_QuickTest_TFG/app/apiREST/controlAccesoProfesor/registro}.
				
		\item \emph{Iniciar sesión en QuickTest:} para ello es necesario incluir en el cuerpo de la petición el email y la constraseña de usuario. La petición es la siguiente: \url{http://localhost/_QuickTest_TFG/app/apiREST/controlAccesoProfesor/login}. 
	\end{list}
	
\end{itemize}

\imagen{ControlAccesoProfesor}{Clase ControlAccesoProfesor}

\paragraph{Clase GestionCuestionario}

La lógica de esta clase será toda aquella que se encuentre relacionada con la gestión sobre los Cuestionarios: \emph{insertar, obtener, duplicar o borrar}. Cabe destacar, que en el controlador del proyecto de partida se ha incluido un método para poder duplicar un cuestionario existente, de esta manera evitamos comunicarnos con el modelo. Todas estas acciones se realizan mediante diferentes peticiones:

\begin{itemize}

	\item \textbf{Get:} esta petición será la encargada de mostrar aquellos cuestionarios dado un identificador de asignatura. \url{http://localhost/_QuickTest_TFG/app/apiREST/gestionCuestionario/id}
	\item \textbf{Post:} esta petición será la encargada de manejar dos lógicas.
		\begin{list}{-}
			\item \emph{Duplicar un cuestionario:} para ello es necesario indicar el identificador del cuestionario. La petición es la siguiente: \url{http://localhost/_QuickTest_TFG/app/apiREST/gestionCuestionario/duplicar/id}.
			\item \emph{Insertar/Editar un cuestionario:} para ello es necesario incluir en el propio cuerpo de la petición si la petición es de insertar o de duplicar. La petición es la siguiente: \url{http://localhost/_QuickTest_TFG/app/apiREST/gestionCuestionario/insertar}.
		\end{list}
	
	\item \textbf{Delete:} esta petición se encarga de eliminar un determinado cuestionario dado un identificador. La petición es la siguiente\url{http://localhost/_QuickTest_TFG/app/apiREST/gestionCuestionario/id}.

\end{itemize}

\imagen{GestionCuestionario}{Clase GestionCuestionario}

Cabe destacar, que estas dos clases explicadas anteriormente no han sido utilizadas finalmente en la aplicación, al no ser necesarias para su correcto funcionamiento. Estas clases fueron desarrolladas durante las primeras fases del proyecto. Estas primeras fases, se dedicaron a la creación del APIREST que posteriormente se utilizaría en la aplicación Android.

\paragraph{Clase SolucionCuestionario}

Esta clase se encarga de la lógica a realizar durante la resolución de un cuestionario: \emph{finalizar o mostrar resultado}.
Todas estas acciones se realizan mediante las siguientes peticiones:

\begin{itemize}

	\item \textbf{Post:} esta petición será la encargada de manejar la lógica a seguir al finalizar un cuestionario. Para ello, obtendrá el contenido del cuestionario resuelto de la petición, almacenará las respuestas del alumno y la nota del mismo. \url{http://localhost/_QuickTest_TFG/app/apiREST/solucionCuestionario/finalizar}.
	
	\item \textbf{Get:} esta petición sera la encargada de obtener la calificación de un alumno que haya resuelto el cuestionario desde la aplicación Android. Para ello, será necesario incluir en la propia petición el identificador del alumno y del cuestionario respectivamente. \url{http://localhost/_QuickTest_TFG/app/apiREST/solucionCuestionario/obtenerNota/idAlumno/idCuestionario}
\end{itemize}

\imagen{SolucionCuestionario}{Clase SolucionCuestionario}

\paragraph{Clase ObtenerCuestionario}

Esta clase se encarga de la obtención de las preguntas y respuestas que forman un cuestionario, así como que preguntas tienen o no comodín y al finalizar un cuestionario, obtener información acerca de como ha sido resuelto el mismo. Destacar,  que en el controlador no existe ningún método que permita obtener un cuestionario por identificador y el estado del mismo. Es por esto que esta funcionalidad ha sido incluida dentro del propio API. Además, ha sido necesario incluir dentro de este la funcionalidad de obtener los comodines verdes, ámbar y la información sobre como ha sido resuelto el cuestionario, ya que los métodos que se encuentran en el controlador corrompen la respuesta al incluir volcados del contenido.

\begin{itemize}

	\item \textbf{Get:} esta petición será la encargada de manejar toda la lógica de la clase.
	\begin{list}{-}
		\item \emph{Obtener cuestionario:} para ello sera necesario incluir en la propia petición el identificador del cuestionario. La petición es la siguiente: \url{http://localhost/_QuickTest_TFG/app/apiREST/obtenerCuestionario/obtener/idCuestionario}.
		
		\item \emph{Obtener el estado de un cuestionario:} es decir, si el cuestionario se encuentra resuelto o no por dicho alumno. Para ello, será necesarios incluir en la propia petición el identificador del alumno y el identificador de cuestionario. La petición es la siguiente:   \url{http://localhost/_QuickTest_TFG/app/apiREST/obtenerCuestionario/estado/idAlumno/idCuestionario}.
		
		\item \emph{Obtener comodín verde}: es decir, permite obtener aquella preguntas que tienen comodín verde, dado un cuestionario. Para ello es necesario incluir en la propia petición el identificador del cuestionario. La petición es la siguiente:\url{http://localhost/_QuickTest_TFG/app/apiREST/obtenerCuestionario/obtenerComodin/verde/idCuestionario}.
		
		\item \emph{Obtener comodín ámbar}: es decir, permite obtener aquella preguntas que tienen comodín ámbar, dado un cuestionario. Para ello es necesario incluir en la propia petición el identificador del cuestionario. La petición es la siguiente:\url{http://localhost/_QuickTest_TFG/app/apiREST/obtenerCuestionario/obtenerComodin/ambar/idCuestionario}.
		
		\item \emph{Obtener información sobre un cuestionario resuelto}: nos permite obtener información sobre somo ha sido resuelto un cuestionario, es decir, en que grupo se encuentra, cuantos comodines a utilizado, etc. Para ello es necesario incluir en la propia petición el identificador de cuestionario y alumno. La petición es la siguiente: \url{http://localhost/_QuickTest_TFG/app/apiREST/obtenerCuestionario/feedback/idCuestionario/idAlumno}.
	
	\end{list}

\end{itemize}

\imagen{ObtenerCuestionario}{Clase ObtenerCuestionario}

\subsubsection{Android}
En la siguiente ilustración podemos ver el diagrama de clases completo de la aplicación Android. Para ello, se ha decidido utilizar el color verde para representar dichas clases.

\begin{landscape}
\imagenAncho{DiagramaClasesAndroid}{Diagrama de clases de la aplicación.}{1.5}
\end{landscape}



\section{Diagrama de paquetes}

Para estos diagramas se ha decido utilizar el mismo criterio de color empleando anteriormente en el diagrama de clases.

\subsection{APIREST}

Este diagrama muestra un desglose de todos los componentes que entran en contacto en el API, como ya se ha comentado en color amarillo corresponde a la lógica dada en el proyecto de partida y en azul el APIREST desarrollado para interactuar con el controlador. Destacar que en color azul representa la estructura que forma el BackEnd de la aplicación. Toda esta estructura se encuentra en  
\_QuickTest\_TFG\_.app.apiREST y esta formada por:

\begin{itemize}

	\item \textbf{modelos:} en este paquete se encuentran todas aquellas clases que se comunican con el controlador del proyecto de partida.
	\item \textbf{test\_apiRest:} en este paquete se encuentran las pruebas unitarias creadas para probar comprobar el correcto funcionamiento del APIREST.
	\item \textbf{utilidades:} en este paquete se encuentran ciertas clases y ficheros de utilidad.
	\item \textbf{vista:} en este paquete se encuentran aquellas clases encargadas de gestionar la lógica de la vista de los datos.

\end{itemize}

\imagen{DiagramPaqRest}{Diagrama de paquetes del APIREST}

\subsection{Aplicación Android}

En este diagrama se muestra el desglose de todos los componentes que forman la arquitectura del FrontEnd de la aplicación. Toda esta estructura se encuentra dentro de QuickTest\_Android.app.src y esta formada por:

\begin{itemize}

	\item \textbf{main.activities:} en este paquete se encuentran todas las actividades necesarias para la construcción de la aplicación.
	\item \textbf{main.adapters:} en este paquete se encuentran todas los adaptadores creados para modificar el layout de las diferentes actividades y fragmentos.
	\item \textbf{main.api:} en esta paquete se encuentran todas clases encargadas de gestionar la correcta conexión de la aplicación con el APIREST y con el web services de Moodle.
	\item \textbf{main.fragments.students:} en este paquete se encuentran todos los fragmentos utilizados para la construcción de la aplicación.
	\item \textbf{main.models:} en este paquete se encuentran todas las clases necesarias para el tratamiento de los datos proporcionados por las diferentes peticiones, tanto al APIREST como al web services de Moodle. Es decir, estas clases son las encargadas de realizar la conversión de los datos en formato JSON a objetos Java o viceversa.
	\item \textbf{main.request:} en este paquete se encuentran todas aquellas clases encargadas de realizar una petición síncrona al APIREST o al web service de Moodle.
	\item \textbf{main.utils:} en este paquete se encuentra una clases encargada de la gestión de las SharedPreferences utilizadas para el correcto funcionamiento del login.
	\item \textbf{test:} en este paquete se encuentran las pruebas unitarias sobre el paquete main.models.
	\item \textbf{androidTest:} en este paquete se encuentra las pruebas automatizadas de interfaz de usuario sobre la aplicación.

\end{itemize}

\imagen{DiagramPaqAndroid}{Diagrama de paquetes de la aplicación Android}

\section{Diseño procedimental}

\subsection{Diagrama de secuencias}

En esta sección se incluyen los diagramas de secuencias o de interacción, encargados de mostrar de una manera mas aproximada la interacción del sistema al realizar cada caso de uso.

\subsubsection{Diagrama de secuencias que implementa CU1: Iniciar sesión}

\imagen{DiagramaSecuenciaCU1}{Diagrama de secuencias caso de uso 1.}

\subsubsection{Diagrama de secuencias que implementa CU2: Olvidar campos}

\imagen{DiagramaSecuenciaCU2}{Diagrama de secuencias caso de uso 2.}

\subsubsection{Diagrama de secuencias que implementa CU3: Recordar campos}

\imagen{DiagramaSecuenciaCU3}{Diagrama de secuencias caso de uso 3.}
\subsubsection{Diagrama de secuencias que implementa CU4: Mostrar cursos}
\imagen{DiagramaSecuenciaCU4}{Diagrama de secuencias caso de uso 4.}

\subsubsection{Diagrama de secuencias que implementa CU5: Mostrar cuestionarios}

\imagen{DiagramaSecuenciaCU5-Teacher}{Diagrama de secuencias caso de uso 5 desde el rol de profesor.}


\imagen{DiagramaSecuenciaCU5-Student}{Diagrama de secuencias caso de uso 5 desde el rol de alumno.}

\subsubsection{Diagrama de secuencias que implementa CU6: Resolver cuestionario}

\imagen{DiagramaSecuenciaCU6}{Diagrama de secuencias caso de uso 6.}

\subsubsection{Diagrama de secuencias que implementa CU7: Mostrar cuestionarios resueltos}

\imagen{DiagramaSecuenciaCU7}{Diagrama de secuencias caso de uso 7.}

\subsubsection{Diagrama de secuencias que implementa CU8: Mostrar calificación de los alumnos}

\imagen{DiagramaSecuenciaCU8}{Diagrama de secuencias caso de uso 8.}
\subsubsection{Diagrama de secuencias que implementa CU9: Ver cuestionarios}
\imagen{DiagramaSecuenciaCU9}{Diagrama de secuencias caso de uso 9.}

\section{Diseño arquitectónico}

A continuación, se muestra el diagrama de despliegue de todo el proyecto, incluyendo la funcionalidad del proyecto de partida. Para ello se utiliza el mismo criterio de color anteriormente utilizado.

\subsubsection{Diagrama de despliegue}

\imagen{DiagramaDespliege}{Diagrama de despliegue del proyecto.}




\apendice{Documentación técnica de programación}

\section{Introducción}

\section{Estructura de directorios}

\section{Manual del programador}

\subsection{Instalación de las herramientas}

\subsubsection{Sistema operativo}

\subsubsection{Servidor}

\begin{itemize}

	\item Acceder a la web oficial de XAMPP en el siguiente enlace:
\url{https://www.apachefriends.org/es/download.html}.

	\item Seleccionar la versión 5.6.30/PHP 5.6.30. Cabe destacar que es necesario instalar esta versión de XAMPP o inferiores, debido a que el proyecto de partida utiliza ciertas funcionalidades que se encuentran sin soporte en versiones superiores. Y comenzamos a descargar el instalador. \imagen{InstXammp}{Instalador de XAMMP}
	
	\item Una vez descargado, procedemos a la instalación de esta herramienta. Para procedemos a realizar la típica instalación de Siguiente – Siguiente y seleccionamos donde ubicar los archivos a instalar.
	
	\item Al finalizar la instalación se lanzará una alerta del firewall de Windows indicando que se necesita conceder acceso a Apache. Marcamos todas las opciones y finalizamos. \imagen{AlertFirApache}{Alerta del firewall de Apache}
	\item Una vez instalado, ya tendremos instalado y disponible nuestro servidor XAMPP. Para el correcto funcionamiento del proyecto de partida y de QuickTest en Android es necesario arrancar Apache y MySQL al iniciar esta herramienta como podemos ver en la siguiente ilustración. \imagen{PanelXammp}{Panel de XAMPP}
	
	\item Además, al arrancar estos servicios nos volverá a saltar el firewall de Windows pidiendo conceder permisos a MySQL. \imagen{AlertFirMysql}{Alerta del firewall de MySQL}
	
	\item Concedemos los permisos y la instalación de nuestro servidor a finalizado.


\end{itemize}

\subsubsection{Moodle}\label{subsec:moodle}

Respecto al proyecto de partida se decidió utilizar como LMS, el mismo utilizado por la Universidad de Burgos, Moodle. Al tener que basarse en este proyecto no se ha realizado ningún cambio en la herramienta utilizada como LMS.
No obstante, en esta ocasión la versión de Moodle empleada es la más actual hasta la fecha Moodle 3.2.2+. Este cambio en la versión utilizada es realizado al tener que utilizar las funciones que proporciona el propio web services de Moodle. De manera que cuanto más actual sea la versión, más funcionalidades tendremos a nuestra disposición.

\begin{itemize}

	\item Descargamos el archivo comprimido de Moodle correspondiente a la versión 3.2.2+: \url{https://sourceforge.net/projects/moodle/files/Moodle/stable32/moodle-3.2.2.zip/download}
	
	\item Seleccionamos la versión comentada anteriormente y comenzamos la descarga del archivo comprimido. \imagen{FolderMoodle}{Carpeta comprimida de Moodle}
	
	\item Una vez finalizada la descarga, descomprimimos el archivo en la carpeta htdocs. Esta carpeta se encuentra en el directorio raíz de XAMPP. Si la ruta de instalación de XAMPP se ha realizado con los campos por defecto, esta se encontrará en ''C:\textbackslash{xampp}\textbackslash{htdocs}''. Esta es la carpeta raíz del localhost, desde donde el servidor XAMPP comenzará a buscar cuando le mandemos una petición.
	
	\item Una vez descomprimido la carpeta abrimos un navegador y realizamos una petición al servidor indicando que queremos acceder a la URL de Moodle. \imagen{urlMoodle}{url de Moodle}
	
	\item Al no tener instalada ni configurada la herramienta, nos pedirá seleccionar el idioma deseado. Una vez seleccionado el idioma, tendremos que seleccionar la ruta de instalación de los nuevos componentes. En este caso dejamos la configuración por defecto como podemos ver en la siguiente ilustración. \imagen{UbicInstDir}{Ubicación de donde instalar los directorios}
	
	\item Posteriormente, tendremos que seleccionar el tipo de base de datos a utilizar. En este caso la opción elegida es MariaDB para evitar problemas de compatibilidad. \imagen{SelBD}{Selección de la base de datos}
	
	\item Al finalizar la selección de la base de datos la instalación nos permitirá modificar la configuración por defecto de la misma. En este caso solamente añadiremos root en la sección Usuario de la base de datos como podemos ver en la siguiente ilustración y aceptamos la licencia. \imagen{ConfMoodle}{Configuración de Moodle}
	
	\item Antes de pulsar en el botón de siguiente tendremos que abrir una nueva pestaña en nuestro navegador e introducir el siguiente enlace localhost/phpmyadmin. \imagen{urlMysql}{url de MySql}
	
	\item Una vez nos encontremos en la página tendremos que cambiar el cotejamiento de la conexión al servidor como podemos ver en la siguiente ilustración.\imagen{CambForBD}{Cambiamos el formato de la base de datos}
	
	\item Una vez realizada la creación de la base de datos de manera manual, podemos volver a la pestaña de instalación de Moodle y proseguir con la instalación.

	\item Si hemos realizado los anteriores pasos de forma correcta el sistema nos comunicará que el entorno cumple todos los requerimientos mínimos y continuamos con la instalación. Este proceso puede durar varios minutos dependiendo de la conexión a Internet.
	
	\item Una vez realizada la instalación de Moodle, es necesario la configuración de la plataforma. Para ello es necesario establecer un administrador de la misma. En nuestro caso la configuración es la que se puede ver en la siguiente ilustración. \imagen{ConfDatAdminMoodle}{Configuramos los datos del administrador de Moodle}
	
	\item A continuación, tendremos que configurar el sitio. Para ello introducimos la siguiente información. \imagen{ConfSites}{Configuramos el sitio}
	
	\item Una vez realizamos de forma correcta todos los anteriores pasos, la instalación y configuración de Moodle ha finalizado

\end{itemize}

\subsubsection{Instalar base de datos}

A continuación, vamos a importar la base de datos que utiliza QuickTest. Esta base de datos es la misma proporcionada por el proyecto de partida. Para llevar a cabo esta tarea tendremos que seguir los siguientes pasos:

\begin{itemize}

	\item Dentro del proyecto de partida nos situamos en la ruta src/InstalarBaseDeDatos.
	
	\item En esta carpeta tendremos un único archivo llamado: instalar\_BaseDatosQuickTest.sql
	
	\item Abrimos una nueva pestaña en nuestro navegador e introducir el siguiente enlace. \imagen{urlMysql}{Url de Mysql}
	
	\item Una vez nos encontremos en el panel de administración de la base de datos tendremos que importar el archivo del punto 2. Para ello seguimos los siguientes pasos. \imagen{FichBD}{Añadimos el fichero de la base de datos}
	
	\item Una vez importado el archivo seleccionamos continuar y el proceso de importación de la base de datos de QuickTest ha finalizado.

\end{itemize}

\subsubsection{Instalar/Configurar QuickTest}

Para instalar QuickTest tendremos copiar la carpeta \_QuickTest\_TFG en la carpeta raíz del localhost de XAMPP para que el servidor pueda buscar el directorio pedido al mandarle una petición.
Para ello tendremos que:

\begin{itemize}

	\item Copiar la carpeta \_QuickTest\_TFG en la carpeta htdocs de XAMPP. \imagen{FolderHtdocs}{Añadimos la carpeta al htdocs}

\end{itemize}

\subsection{Registrar un usuario en QuickTest}

Para poder utilizar QuickTest desde un LMS, en este caso Moodle, tendremos que estar previamente registrado en QuickTest. Para ello tendremos que:

\begin{itemize}

	\item Abrir una nueva pestaña en el navegador e introducir el siguiente enlace: \imagen{UrlIniQT}{Abrimos la página de inicio de QuickTest}
	
	\item Seleccionamos Acceder/Registrarse. Una vez seleccionado este botón aparecerá una ventana emergente que permitirá Entrar o Registrarse.
	
	\item Para llevar a cabo el registro de un nuevo usuario en QuickTest tendremos que seguir dos criterios importantes.
	\begin{itemize}
         \item En el campo Email tendremos que usar el mismo que en Moodle.
         \item En el campo Contraseña podremos utilizar cualquiera siempre y cuando la contraseña este formada entre 5 y 10 caracteres.
    \end{itemize}
    En la siguiente ilustración podemos observar un inicio de sesión. Este usuario se encuentra por defecto en la base de datos inicial por lo que si se desea podemos saltarnos el paso de registrar un usuario. La contraseña es 12345. \imagen{EntQT}{Entramos en QuickTest}
    \item Una vez realizado el registro y su posterior inicio de sesión el sistema informará con una ventana de bienvenida. \imagen{PanFF}{Panel de primeros pasos}
    
    \item Al pinchar sobre Publicar cuestionario se abrirá una ventana emergente. Es en esta ventana donde se encuentra los tres campos necesarios para poder publicar y enlazar nuestros cuestionarios con Moodle. La información necesarios es la siguiente. \imagen{WPubCues}{Ventana emergente de publicar un cuestionario}
    
	
\end{itemize}


\subsection{Configurar/Publicar un cuestionario en Moodle}

En esta sección vamos a configurar Moodle para poder utilizar QuickTest. Para llevar a cabo esta tarea tendremos que seguir los siguientes pasos.

\subsubsection{Configuración}


\begin{itemize}

	\item Abrimos una nueva pestaña del navegador y accedemos a Moodle.
	
	\item Una vez allí tendremos que iniciar sesión con el nombre de usuario y contraseña propios del administrador del sistema, en este caso son las credenciales utilizadas en la sección \ref{subsec:moodle}. \imagen{InitSesMoodle}{Inicio de sesión de Moodle}
	
	\item Una vez iniciada la sesión tendremos que crear un curso donde añadir nuestro cuestionario, para ello tendremos que:
	
	\begin{itemize}
	
		\item Entramos en Administración del sitio.
		
		\item Una vez allí entramos en la categoría de Cursos y dentro de esa sección en Administrar cursos y categorías. \imagen{CrearCurso}{Creamos un curso}
		
		\item Una vez en esa sección seleccionamos Crear un nuevo curso y rellenamos los campos con la siguiente información. Al finalizar guardamos la información.
		\imagen{RelCampos}{Rellenamos los campos del nuevo curso}
	\end{itemize}
	
	\item Al crear el curso tendremos que matricular a los usuarios. Para ello vamos a crear dos usuarios.
	
	\begin{itemize}
		
		\item Tendremos que entrar en Administración del sitio, una vez allí tendremos que ir a la categoría de Usuarios y dentro de allí a la sección Crear un nuevo usuario. \imagen{CreNueUsers}{Creamos los nuevos usuarios}
		
		\item A continuación, tendremos que crear dos usuarios como podemos ver en las siguientes ilustraciones. \imagen{RelCmps}{Rellenamos los campos del primer usuario} \imagen{RelCmps2}{Rellenamos los campos del segundo usuario}
	\end{itemize}
	
	\item Una vez los usuarios tendremos que matricular dichos usuarios en el curso creado para ello tendremos que entrar en Administración del sitio, una vez allí tendremos que ir a la categoría de Cursos y dentro de allí a la sección Administrar cursos y categorías.
	
	\item Allí seleccionamos el nuevo curso creado y pinchamos en Usuarios matriculados.
	
	\item A continuación, matriculamos ambos usuarios creados anteriormente, el alumno con rol de Estudiante y el profesor con el rol de Profesor.
	
	

\end{itemize}

\subsubsection{Publicar}

Para llevar a cabo la publicación de un cuestionario tendremos que:

\begin{itemize}

	\item Tendremos que entrar en Administración del sitio, una vez allí tendremos que ir a la categoría de Cursos y dentro de allí a la sección Administrar cursos y categorías.
	
	\item Una vez allí, seleccionamos el curso creado y en Vista.
	
	\item A continuación, activamos la edición y añadimos una herramienta externa.
	
	\begin{itemize}
	
			\item Para ello, tendremos que seleccionar Agregue una actividad o recurso. \imagen{AgHE}{Agregamos una nueva herramienta externa}
			\item En este instante comenzamos a configurar la herramienta externa para poder utilizar QuickTest. Tendremos que rellenar los campos con la siguiente información.\imagen{ConfHE2}{Configuramos la nueva herramienta externa 1} \imagen{ConfHE}{Configuramos la nueva herramienta externa 2}
			
			\item Además, revisar que en el apartado Privacidad se encuentren activas las siguientes opciones. \imagen{CamPrivi}{Campos de la privacidad}
	
	\end{itemize}
	
	\item En este instante la publicación ha finalizado.

\end{itemize}







\section{Compilación, instalación y ejecución del proyecto}

\section{Pruebas del sistema}

\apendice{Documentación de usuario}

\section{Introducción}

Este manual se presentan y explican los requisitos necesarios para utilizar la aplicación, como instalar en el un dispositivo físico y otras pautas de configuración para su correcto funcionamiento.

\section{Requisitos de usuarios}

Para poder utilizar la aplicación es necesario cumplir una serie de requisitos mínimos:

\begin{itemize}

	\item Tener instalado  las herramientas explicadas en \ref{subsec:RequisitosMinimos}.
	\item Configurar las herramientas como se encuentran explicadas en \ref{subsubsec:InstHerramientas}.

\end{itemize}

\section{Instalación}

\subsection{Configuración}\label{subsec:ConfiguracionMaquina}

La aplicación ha sido desarrollada para dispositivos Android 4.1 (Jelly Bean) o superiores. Antes de iniciar la aplicación es necesario realizar unos pequeños cambios en la configuración.

\begin{enumerate}
	
	\item Modificar la dirección ip predeterminada del archivo de configuración de Moodle. Para ello tendremos que:
	\begin{enumerate}
		\item Situarnos dentro de la propia carpeta de la herramienta: 

D:\textbackslash{}xampp\textbackslash{}htdocs\textbackslash{}moodle
		\item Una vez situados en dicha ruta tendremos que modificar la dirección del host: \emph{localhost} por la dirección ip de la propia máquina del archivo config.php como podemos ver en la siguiente ilustración.
		\imagen{archivoConfig}{Archivo config.php de Moodle.}
	\end{enumerate}		
	
	\item Al haber modificado el host de Moodle tendremos que modificar la dirección de inicio de todas nuestras herramientas externas QuickTest por la dirección ip de la maquina. Esto es porque tenemos QuickTest y Moodle en la misma máquina, sino no seria necesario. \imagenAncho{modUrlInicio}{Modificando la URL de inicio de las herramientas externas.}{1}
	
	\item A continuación, tendremos que abrir nuestro proyecto en Android Studio y modificar los siguientes archivos dentro del paquete api. \imagenAncho{ModificarapiAndroid}{Archivos a modificar dentro del paquete api.}{0.8}
	
	\begin{enumerate}
		
		\item \emph{APIMoodle:} tendremos que modificar la variable \emph{BASE} por la dirección de la máquina y las variables \emph{USERNAME, PASSWORD} por el nombre y contraseña del usuario encargado de manejar el web services de Moodle, ya que el es el único con privilegios para realizar las peticiones de información. \imagen{ModAPIMoodle}{Modificamos las variables de la clase APIMoodle}
		
		\item \emph{APIRest:} tendremos que modificar la variable \emph{BASE} por la dirección de la máquina. \imagen{ModAPIRest}{Modificamos la variable de la clase APIRest.}
		
	\end{enumerate}
	 
\end{enumerate}

Una vez realizada esta breve modificación en la configuración, la aplicación Android se encuentra completamente operativa para su ejecución desde el propio emulador de Android Studio o desde nuestro propio dispositivo Android.

\subsection{Instalación en un dispositivo Android físico}

Para llevar cabo la instalación de la aplicación en nuestro propio dispositivo Android es necesario:

\begin{itemize}
	\item Haber realizado la breve configuración previa del punto anterior. \ref{subsec:ConfiguracionMaquina}.
	
	\item Conectar nuestro dispositivo al ordenador.
	
	\item Nuestro dispositivo pertenezca a la misma red que nuestro equipo.
\end{itemize}

Una vez realizado estos pasos tendremos que:

\begin{enumerate}
	\item Arrancar el proyecto en Android Studio.
	\item Ejecutar la aplicación y seleccionar como dispositivo de destino, nuestro móvil.\imagen{InstApp}{Instalación de la aplicación en un dispositivo físico.}
\end{enumerate}

Una vez realizados estos pasos podremos utilizar nuestra aplicación desde nuestro propio dispositivo Android.

\section{Manual del usuario}

En esta sección se describen el funcionamiento de las funcionalidades desarrolladas de la aplicación.

\subsection{Iniciar Sesión}

Al iniciar la aplicación por primera vez nos aparecerá la ventana de \emph{inicio de sesión.} Una vez allí tendremos que:

\begin{enumerate}
	\item Introducir nuestro nombre de usuario y contraseña de Moodle.
	\item Si se quiere recordar estas credenciales para que en futuros acceso esta ventana no se active, bastará con activar la opción \emph{Recordar campos.}
	\item Click en \emph{Iniciar Sesión.}
\end{enumerate}

Una vez iniciada sesión la aplicación nos mostrará los cursos en los que el usuario se encuentra matriculado independientemente de su rol en los mismos.

\imagen{IniciarSesion}{Iniciar sesión en la aplicación.}



\subsection{Obtener la calificación de los alumnos}

Para poder comprobar la calificación de los alumnos de un profesor simplemente tendremos que:

\begin{enumerate}
	\item Entrar en el curso.
	\item Seleccionar el cuestionario a comprobar.
	\item Comprobar la calificación de cada alumno.
\end{enumerate}

\imagen{ObtenerCalAlumnos}{Obtener la calificación de los alumnos matriculados en un curso.}

Como podemos observar en la anterior ilustración, se utiliza el siguiente criterio para determinar que cuestionarios han sido resueltos de la aplicación web, desde la aplicación Android o aun no han sido resueltos.

\begin{itemize}
	\item \emph{Icono de Moodle:} cuestionario resuelto desde la aplicación web.
	
	\item \emph{Icono de información:} cuestionario sin resolver.
	
	\item \emph{Calificación:} cuestionario resuelto desde la aplicación Android.
\end{itemize}

\subsection{Mostrar cuestionario}

Para obtener el cuestionario ya completado por los alumnos tendremos que:

\begin{enumerate}
	\item Entrar en el curso.
	\item Seleccionar el cuestionario a comprobar.
	\item Seleccionar, \emph{Ver cuestionario}.
\end{enumerate}

\imagen{VerCuestionario}{Obtenemos el cuestionario al que se han enfrentado los alumnos.}

\subsection{Resolver un cuestionario}

Para resolver un cuestionario tendremos que:
\begin{enumerate}
	\item Entrar en el curso.
	\item Posicionarnos en \emph{Mis cuestionarios sin resolver} en el \emph{Burger Menu}.
	\item Click largo sobre el cuestionario a resolver y seleccionamos \emph{Resolver}.
	\item Una vez resueltas las preguntas, click en \emph{Enviar y finalizar cuestionario} y en la ventana emergente seleccionar  \emph{Enviar.}
\end{enumerate}

\imagen{ResolverCuestionario}{Resolvemos un cuestionario.}


\subsection{Revisar un cuestionario}

Para revisar un cuestionario resuelto tendremos que:
\begin{enumerate}
	\item Entrar en el curso.
	\item Posicionarnos en \emph{Mis cuestionarios resueltos} en el \emph{Burger Menu.}
	\item Click largo sobre el cuestionario a revisar y seleccionamos \emph{Revisar.}
	\item Si deseamos obtener información adicional sobre como ha sido resuelto el cuestionario, hacer click en \emph{Ver resultados}
\end{enumerate}

\imagen{RevisarCuestionario}{Revisar un cuestionario.}

\subsection{Cerrar sesión}

Para cerrar sesión u olvidar y cerrar sesión tendremos que:

\begin{itemize}
	\item \emph{Rol alumno:}
	\begin{enumerate}
		\item Entrar en el curso.
		\item Abrir en \emph{Burger Menu} y seleccionar \emph{Cerrar sesión u Olvidar y cerrar sesión.}
	\end{enumerate}
	\item \emph{Rol Profesor:}
		\begin{enumerate}
			\item Entrar en el curso.
			\item Abrir el botón de ajustes y seleccionar \emph{Cerrar sesión u Olvidar y cerrar sesión.}
		\end{enumerate}
\end{itemize}

\imagen{CerrarSesion}{Cerrar sesión desde ambos roles.}



\bibliographystyle{plain}
\bibliography{bibliografiaAnexos}

\end{document}
